\documentclass[11pt,a4paper]{article}

\usepackage[T1]{fontenc}
\usepackage[utf8]{inputenc}
\usepackage{lmodern}
\usepackage[indonesian]{babel}
\usepackage[a4paper,margin=2.5cm]{geometry}
\usepackage{setspace}
\usepackage{enumitem}
\usepackage{parskip}
\usepackage[hidelinks]{hyperref}

\setstretch{1.3}
\setlist[enumerate]{itemsep=7pt, topsep=4pt, leftmargin=*}

\begin{document}

\begin{center}
  {\large\bfseries Daftar Pertanyaan Wawancara}\\[3pt]
  {\normalsize Pengalaman Berjalan Kaki di Koridor Sudirman, Jakarta}\\[8pt]
  \rule{\textwidth}{0.6pt}
\end{center}

\vspace{6pt}

Terima kasih atas kesediaan Bapak menerima kami. Kami Tim Riset Pijak, sekelompok anak muda yang
sedang belajar di Apple Developer Academy @ BINUS, dan kami sedang mempelajari pengalaman berjalan
kaki di koridor Sudirman, terutama bagi warga yang setiap hari bekerja dan mengandalkan
transportasi umum di sana. Lewat wawancara sekitar 60 menit ini, kami ingin belajar dari cara
Bapak memandang dan menata ruang bagi pejalan kaki. Pertanyaan berikut hanya garis besar; kami
justru senang bila perbincangan berkembang ke arah yang Bapak rasa penting.

\section*{Jakarta dari masa ke masa}
\begin{enumerate}
  \item Bagaimana Bapak melihat perubahan cara Jakarta memperlakukan pejalan kaki dari waktu ke
        waktu?
  \item Menurut Bapak, kota yang ideal itu berangkat dari filosofi seperti apa?
\end{enumerate}

\section*{Cerita koridor Sudirman}
\begin{enumerate}[start=3]
  \item Sebagai pendatang yang sehari-hari naik transportasi umum, kami merasa Sudirman nyaman
        sekali untuk berjalan kaki. Boleh Bapak ceritakan pendekatan di balik penataannya?
  \item Waktu menata Sudirman, bagaimana Bapak melihat orang-orang yang berjalan di sana, terutama
        para pekerja dan pengguna transportasi umum yang lalu-lalang setiap hari?
  \item Boleh Bapak ceritakan prosesnya? Apa saja yang harus Bapak timbang dan seimbangkan,
        misalnya fungsi, kapasitas, kenyamanan, dan keindahan?
\end{enumerate}

\section*{Ketertiban dan ruang untuk semua}
\begin{enumerate}[start=6]
  \item Bagaimana Bapak melihat peran ketertiban dan perawatan bersama dalam menjaga ruang publik
        tetap nyaman?
  \item Bagaimana caranya membuat ruang yang nyaman untuk semua orang, termasuk lansia dan
        penyandang disabilitas, seperti yang Bapak upayakan di Sudirman?
\end{enumerate}

\section*{Data dan teknologi}
\begin{enumerate}[start=8]
  \item JAKI lahir di masa kepemimpinan Bapak. Boleh Bapak ceritakan bagaimana ide menghubungkan
        warga dengan pemerintah lewat teknologi itu tumbuh, dan apa yang Bapak pelajari darinya?
\end{enumerate}

\section*{Refleksi dan keberlanjutan}
\begin{enumerate}[start=9]
  \item Dari semua pengalaman menata kota, pelajaran apa yang paling Bapak bawa sampai sekarang?
  \item Menurut Bapak, apa yang membuat sebuah penataan kota bisa bertahan dan terus terawat dalam
        jangka panjang?
\end{enumerate}

\section*{Penutup}
Apa pesan dan harapan Bapak untuk anak-anak muda yang ingin ikut membuat kota yang lebih baik bagi
pejalan kaki?

\vspace{8pt}
\noindent\rule{\textwidth}{0.3pt}

\vspace{2pt}
{\small Hasil wawancara ini kami pakai untuk belajar, bukan untuk kepentingan komersial. Kami
merekam hanya bila Bapak mengizinkan, dan kami akan mengirimkan catatannya untuk Bapak periksa
sebelum kami gunakan.}

\end{document}
